% Three-fold out-of-fold feature construction for Stacking.
% Schematic only; no empirical quantities are encoded.
% Canvas: 164 mm x 68 mm. Dependencies: project palette and TikZ.
\begin{tikzpicture}[
  x=1mm,y=1mm,
  font=\figurefont\footnotesize,
  text=BookInk,
  fold/.style={
    rectangle,
    model input,
    minimum width=16mm,
    minimum height=9mm,
    align=center,
    inner sep=1mm
  },
  train/.style={
    rectangle,
    model context,
    minimum width=27mm,
    minimum height=9mm,
    align=center,
    inner sep=1mm
  },
  model/.style={
    rectangle,
    model hidden,
    minimum width=25mm,
    minimum height=9mm,
    align=center,
    inner sep=1mm
  },
  predict/.style={
    rectangle,
    model training,
    minimum width=25mm,
    minimum height=9mm,
    align=center,
    inner sep=1mm
  },
  matrix/.style={
    rectangle,
    model output,
    minimum width=18mm,
    minimum height=43mm,
    text width=14mm,
    align=center,
    inner sep=1mm
  },
  meta/.style={
    rectangle,
    model output,
    minimum width=24mm,
    minimum height=9mm,
    align=center,
    inner sep=1mm
  },
  flow/.style={
    model flow
  }
]
  \path[use as bounding box] (0,0) rectangle (164,68);

  \node[anchor=west,font=\figurefont\bfseries\small] at (1,64) {三折轮换};

  \node[fold] (v1) at (12,52) {留出$V_1$};
  \node[fold] (v2) at (12,35) {留出$V_2$};
  \node[fold] (v3) at (12,18) {留出$V_3$};

  \node[train] (t1) at (42,52) {训练$D_n\setminus V_1$};
  \node[train] (t2) at (42,35) {训练$D_n\setminus V_2$};
  \node[train] (t3) at (42,18) {训练$D_n\setminus V_3$};

  \node[model] (m1) at (75,52) {$h_1^{(-1)},\ldots,h_M^{(-1)}$};
  \node[model] (m2) at (75,35) {$h_1^{(-2)},\ldots,h_M^{(-2)}$};
  \node[model] (m3) at (75,18) {$h_1^{(-3)},\ldots,h_M^{(-3)}$};

  \node[predict] (p1) at (108,52) {预测$V_1$};
  \node[predict] (p2) at (108,35) {预测$V_2$};
  \node[predict] (p3) at (108,18) {预测$V_3$};

  \node[matrix] (z) at (145,35)
    {按原样本顺序拼接\\[1mm]
     折外矩阵\\$\bm{Z}\in\mathbb{R}^{n\times M}$};
  \node[meta] (g) at (145,5.5) {元学习器$g$};

  \draw[flow] (v1) -- (t1);
  \draw[flow] (t1) -- (m1);
  \draw[flow] (v2) -- (t2);
  \draw[flow] (t2) -- (m2);
  \draw[flow] (v3) -- (t3);
  \draw[flow] (t3) -- (m3);
  \draw[flow] (m1) -- (p1);
  \draw[flow] (m2) -- (p2);
  \draw[flow] (m3) -- (p3);
  \draw[flow] (p1.east) -- ([yshift=17mm]z.west);
  \draw[flow] (p2.east) -- (z.west);
  \draw[flow] (p3.east) -- ([yshift=-17mm]z.west);
  \draw[flow] (z.south) -- (g.north);
\end{tikzpicture}
